% Generated draft: requires mathematical and source review.
\documentclass[11pt]{article}
\usepackage[margin=25mm]{geometry}
\usepackage{fontspec}
\setmainfont{Latin Modern Roman}
\usepackage{amsmath,amssymb,mathtools,unicode-math}
\setmathfont{Latin Modern Math}
\usepackage{enumitem,needspace,xurl,hyperref}
\setlength{\parindent}{0pt}
\setlength{\parskip}{5pt}
\setlength{\emergencystretch}{3em}
\newcommand{\N}{\mathbb{N}}
\newcommand{\Z}{\mathbb{Z}}
\newcommand{\Q}{\mathbb{Q}}
\newcommand{\R}{\mathbb{R}}
\newcommand{\C}{\mathbb{C}}
\newcommand{\F}{\mathbb{F}}
\newcommand{\A}{\mathbb{A}}
\newcommand{\PP}{\mathbb{P}}
\newcommand{\E}{\mathbb{E}}
\newcommand{\eps}{\varepsilon}
\newcommand{\dd}{\,\mathrm{d}}
\begin{document}
\section*{UnsolvedMath ID: 30005691}
\textbf{Source problem code:} catalogue entry
\subsection{Definitions and mathematical statement}
A rough path is an enhanced driving signal equipped with iterated-integral data satisfying algebraic compatibility and analytic regularity conditions. The precise enhancement, algebraic structure, regularity parameters, and integration theory depend on the class under consideration; the supplied record does not fix them. Throughout the requested problem, integrals are to be interpreted in the rough sense. An integration-by-parts identity means a product identity for the relevant rough integrals, including any correction terms required by the enhancement. An It\^o-type formula means a change-of-variable identity with the corresponding correction terms. The source does not supply either identity explicitly.

The research question is: which additional classes of rough paths admit integration-by-parts identities from which an It\^o-type change-of-variable formula can be derived? Here ``additional'' refers to classes beyond those discussed in the source context, which includes branched and quasi-geometric rough paths; the supplied dated assessment also mentions planarly branched results. For a precise assertion about any proposed class, the admissible rough paths, integrands, test functions, drift terms, regularity assumptions, and identities must be specified. These data are absent from the supplied excerpt, so this question cannot presently be expressed as a fully quantified theorem or conjecture.

The supplied excerpt gives the following qualifications to its existing results. In the quadratic-drift case, described there as an It\^o stochastic differential equation, its formula~(1) coincides with the classical It\^o formula. The definition of an object written \(Fq\) is said to rely entirely on the pre-Lie structure of vector fields on Euclidean space. The proof is said to use an extended Davie expansion with drift-induced corrections. Quasi-geometric rough paths are said to correspond to a particular quotient of branched rough paths, yielding an analogue of formula~(1), connected to the excerpt's references~[1, 5]. Formula~(1), the definition and intended typesetting of \(Fq\), the quotient, and the identities of references~[1, 5] are not supplied. These assertions are preserved as source context, not independently verified mathematical claims. In particular, no unrestricted implication from an unspecified integration-by-parts identity to an It\^o formula is asserted here.
\subsection{Short English statement}
Identify further classes of rough paths for which rough integration admits a product rule with suitable correction terms and a corresponding It\^o change-of-variable formula. The supplied excerpt omits the formulas and analytic hypotheses needed for a fully precise formulation. Its partially solved label and descriptions of existing results are supplied curation requiring independent review.
\subsection{Sources}
\begin{itemize}
\item Supplied local UnsolvedMath registry record, permanent ID \(30005691\), source problem code \texttt{catalogue entry}, title ``Integration by Parts and It\^o Formulas for Rough Paths.'' Project registry locations are \texttt{data/unsolvedmath/problems.json} and \texttt{catalog/problems.json}; the supplied record is the evidence used here. No individual proposer is identified. The proposed year is recorded as 2023.
\item Supplied report citation: \emph{Mini-Workshop: Combinatorial and Algebraic Structures in Rough Analysis and Related Fields (2024)}, Oberwolfach Reports, pp.~3063--3102, DOI \texttt{10.4171/owr/2023/54}. \url{https://doi.org/10.4171/owr/2023/54}. The extracted passage includes the page marker 3095. The citation's 2024 designation and its 2023 DOI component and proposed year are retained without resolving their relationship.
\item Supplied literature source, labelled ``Bellingeri-Ferrucci-Tapia'': \url{https://arxiv.org/abs/2312.04523}. No title or further bibliographic details are supplied.
\item Supplied literature source, labelled ``Li-Gao'': \url{https://arxiv.org/abs/2501.11886}. No title or further bibliographic details are supplied.
\item Literature assessment dated exactly \texttt{2026-08-22}: the supplied curation labels the problem partially solved and states that integration-by-parts and It\^o-type formulas are established for branched and quasi-geometric rough paths, with newer planarly branched results, while further classes remain an open direction in the report. This is supplied dated curation, not an independently verified assessment; its status and cited claims require independent verification.
\item The registry labels the statement \texttt{corrected\_verified} and reports that the official report was checked directly and extraction artifacts removed. This verification claim belongs to the supplied record; no source visit or independent verification is claimed in this draft. The excerpt's references~[1, 5] cannot be identified from the supplied material.
\end{itemize}
\end{document}
